% !TeX root = ../traffic_rules.tex

\section{Introduction}
This document serves as an overview about formalized traffic rules based on the German Road Traffic
Regulation (StVO), the Vienna Convention on Road Traffic \cite{vienna_convention} (VCoRT), judicial decisions, and comments by lawyers. 

The formalization process of traffic rules in temporal logic can be separated into four steps:
\begin{enumerate}
	\item \textbf{Extracting rules from legal sources}: Based on a traffic rule, all related rules can be extracted from legal sources. Legal sources can contain judicial decisions, other traffic law sections, and consultancy of layers.  
	
	\item \textbf{Concretizing extracted rules}: The legal texts and judicial decisions are combined and concretized using natural language. As part of the concretization, the situation when the rule is applicable should be expressed. It can also be the case that a single traffic rule addresses different aspects and therefore each has to be considered separately, e.g., \stvo$\S 4(1)$ addresses the distance between vehicles and also restricts the braking of a vehicle. The use of natural language allows lawyers to evaluate the concretization so that consistency between the formalized rules for autonomous vehicles and the traffic rules for human traffic participants can be ensured.
	
	\item \textbf{Extracting functions, predicates, and propositions}: 
	We use predicates, functions, and atomic propositions for the evaluation of traffic situations, e.g., the allowed speed limit or the velocity of a preceding vehicle. These elements should be defined in a modular way such that they can be easily reused.
	
	\item \textbf{Creating temporal logic formulas}: The extracted predicates, functions, and propositions have to be combined to a temporal logic formula matching the concretized sentence from 2).
\end{enumerate}

The remainder of this paper is structured as follows. Sec.~\ref{sec:preliminaries} introduces mathematical foundations. Sec.~\ref{sec:formalization} presents the formalization of traffic rules. In Sec.~\ref{sec:predicates} predicates and auxiliary functions for the temporal logic formulas are defined.